\uuid{k28s}
\exo7id{1950}
\titre{exo7 1950}
\auteur{legall}
\organisation{exo7}
\datecreate{2003-10-01}
\isIndication{false}
\isCorrection{false}
\chapitre{Série de Fourier}
\sousChapitre{Autre}
\cpge{mpsi-mp2i.calcul-algebrique-trigonometrie.d, pcsi.calcul-algebrique-trigonometrie.d, ptsi.calcul-algebrique-trigonometrie.d, mpsi-mp2i.nombres-complexes.c, pcsi.nombres-complexes.c, ptsi.nombres-complexes.c, mpsi-mp2i.integration.c, pcsi.integration.b, ptsi.integration.b, bcpst1.trigonometrie, bcpst1.methodes-calcul, mp-mpi.suites-series-fonctions.c}
\module{Analyse}
\niveau{L2}
\difficulte{}

\contenu{
\texte{
Soit $f$ une fonction int\'egrable au sens de Riemann p\'eriodique de 
p\'eriode $2\pi $. On d\'esigne par~:
$\displaystyle{ \frac{a_0}{2}+ \sum _{k=1}^\infty(a_k \cos (kx)+ b_k 
\sin (kx))}$ sa s\'erie de Fourier et on pose, pour tout $n\in \Nn :
\displaystyle{ S_n(x)=\frac{a_0}{2}+ \sum _{k=1}^n(a_k \cos (kx)+ b_k 
\sin (kx))}.$
}
\begin{enumerate}
    \item \question{Soit $\theta \in \Rr-2\pi \Zz .$ Montrer $\displaystyle{ 
\frac{1}{2}+ \sum _{k=1}^n \cos (k\theta)=\frac{\sin 
(n+\frac{1}{2})\theta}
{2\sin \frac{\theta }{2}}}.$}
    \item \question{Etablir que $
\displaystyle{ S_n(x)=\frac{1}{\pi}\int _{-\pi}^\pi \frac {\sin 
(n+\frac{1}{2})(x-t)}{2\sin \frac{(x-t)}{2}}f(t)dt}$.}
    \item \question{En d\'eduire $ \displaystyle{ S_n(x)=\frac{1}{2\pi}\int 
_{-\pi}^\pi f(x+\theta ) \frac{\sin (n+\frac{1}{2})\theta}
{\sin \frac{\theta }{2}}d\theta } $.}
    \item \question{Calculer $\displaystyle{ \int _{0}^\pi \frac{\sin (n+\frac{1}{2})\theta}
{\sin \frac{\theta }{2}}d\theta } .$}
\end{enumerate}
}
